\section{Model Performans Karşılaştırması}

Dört farklı makine öğrenmesi algoritması ile eğitilen modellerin test seti üzerindeki performansları Tablo \ref{tab:sonuclar}'da özetlenmiştir.

\begin{table}[H]
\centering
\caption{Model Performans Metrikleri}
\label{tab:sonuclar}
\begin{tabular}{lrrrrr}
\toprule
Model & Accuracy & Precision & Recall & F1-Score & ROC-AUC \\
\midrule
Lojistik Regresyon & - & - & - & - & - \\
Random Forest & - & - & - & - & - \\
XGBoost & - & - & - & - & - \\
LightGBM & - & - & - & - & - \\
\bottomrule
\end{tabular}
\end{table}

\textit{Not: Yukarıdaki tablo, notebook'lar çalıştırıldıktan sonra gerçek değerlerle güncellenmelidir.}

\section{ROC Eğrileri}

Şekil \ref{fig:roc} tüm modellerin ROC eğrilerini göstermektedir.

\begin{figure}[H]
\centering
\includegraphics[width=0.8\textwidth]{roc_curves.png}
\caption{Modellerin ROC Eğrileri}
\label{fig:roc}
\end{figure}

\section{Confusion Matrix Analizi}

Her model için confusion matrix Şekil \ref{fig:cm}'de gösterilmektedir.

\begin{figure}[H]
\centering
\includegraphics[width=0.9\textwidth]{confusion_matrices.png}
\caption{Model Confusion Matrix'leri}
\label{fig:cm}
\end{figure}

\section{Özellik Önem Analizi}

Random Forest modeli için en önemli 15 özellik Şekil \ref{fig:importance}'de gösterilmektedir.

\begin{figure}[H]
\centering
\includegraphics[width=0.8\textwidth]{feature_importance.png}
\caption{En Önemli 15 Özellik (Random Forest)}
\label{fig:importance}
\end{figure}

Analiz sonuçlarına göre en önemli özellikler:

\begin{enumerate}
    \item Önceki günlerin yıldırım durumu (lag özellikleri)
    \item Mevsimsel faktörler
    \item Rakım
    \item Rolling window istatistikleri
\end{enumerate}

\section{Çapraz Doğrulama Sonuçları}

En iyi performans gösteren model için 5-katlı stratifiye çapraz doğrulama uygulanmıştır:

\begin{itemize}
    \item Ortalama F1-Score: [değer]
    \item Standart Sapma: [değer]
\end{itemize}

Bu sonuçlar, modelin farklı veri bölümlerinde tutarlı performans gösterdiğini ortaya koymaktadır.

\section{Precision-Recall Analizi}

Dengesiz veri setleri için önemli olan Precision-Recall eğrileri Şekil \ref{fig:pr}'de gösterilmektedir.

\begin{figure}[H]
\centering
\includegraphics[width=0.8\textwidth]{precision_recall_curves.png}
\caption{Precision-Recall Eğrileri}
\label{fig:pr}
\end{figure}
